\documentclass[11pt]{article}
\usepackage{manual_docs_style}
\externaldocument{build/candidate_generation}[candidate_generation.pdf]
\externaldocument{build/dataset_manifest_schema}[dataset_manifest_schema.pdf]

\begin{document}

\docTitle{Web Applications}
\phantomsection\label{docstart:web_model_explorer}
This document describes the two web applications available in the repo.
\section*{Model Explorer}

This is a local read-only web explorer used to inspect materialized \name{} runs and generated question bundles.
By default, the data root is the repository's \code{models/} directory, which should contain \code{simulink/<simulator_id>/}.
For an external bundle, set:
\begin{DocCode}
WEB_MODEL_EXPLORER_MODEL_ARTIFACT_DIR=/absolute/path/to/simulator
\end{DocCode}
The value is the simulator's artifact directory itself.

Start it from the repository with:

\begin{DocCode}
bash web_model_explorer/start.sh
\end{DocCode}

The launcher starts the Next.js app in \code{web_model_explorer/}.
It uses \code{0.0.0.0:3001} by default.
The following variables override the host, port, and simulator artifact directory, respectively:
\begin{DocCode}
WEB_MODEL_EXPLORER_HOST
WEB_MODEL_EXPLORER_PORT
WEB_MODEL_EXPLORER_MODEL_ARTIFACT_DIR
WEB_MODEL_EXPLORER_EVALUATION_ROOT
WEB_MODEL_EXPLORER_EVALUATION_TIME0_ROOT
WEB_MODEL_EXPLORER_EVALUATION_ALTERNATIVE_PROBE_ROOT
WEB_MODEL_EXPLORER_PAPER_EVALUATION_MANIFEST
\end{DocCode}
The artifact-directory override is a direct path to one simulator, not a root containing \code{simulink/}.
The evaluation-root override points to the released bundle containing the
\code{BallDrop/}, \code{BounceBall/}, and \code{MassSlide/} directories.
The evaluation-time-zero-root override points to a separate, viewer-only
auxiliary artifact containing regenerated clean time-zero counterparts.
The evaluation-alternative-probe-root override points to the separate,
viewer-only artifact containing optimized clean alternative-root-cause probes.
The paper-evaluation-manifest override points to the authoritative
\code{balanced_heldout_evaluation.json} artifact.
Run-data and noise helpers use a Python interpreter.
The current application accepts \code{WEB_MODEL_EXPLORER_PYTHON} (or \code{PYTHON_BIN}) as an explicit interpreter override before checking repository virtual-environment paths and then \code{python3}.
It should have a structure such as:
\begin{DocCode}
    experiment_config.json
    description_levels.json
    generated/metadata.json
    plans/<policy_id>/
    runs/model_record.json
    runs/<run_id>/
\end{DocCode}




The registry is loaded with a paged approach to handle large policies (i.e. the design principle is that we fetch only what is needed).
The first page contains N families (by default N=250), and additional pages are fetched on demand (by clicking the next-page button).
The \code{Eligible only} checkbox is unchecked by default.
When selected, it reloads the first page and shows only families containing at least one intervention run whose \code{cheap_filter_pass} value is explicitly \code{true}; expanding a retained family likewise shows only its passing intervention runs.
When the metrics file contains explicit family summaries, their \code{family_cheap_filter_pass} value determines whether a family is retained.
For run-only metric files, the explorer derives the family result as ``at least one passing run.''
Families whose runs all failed the cheap filter and families without a known eligibility result are excluded.
Filtering is applied before pagination, so the displayed family count and page controls describe the complete eligible set rather than only the eligible rows from the previously visible page.
Full family details, including intervention and probe runs, are loaded when a family is expanded (i.e. when a family is clicked).
Deep links are supported and can open a selected run without requiring the full policy registry to be loaded up front.

For large resolved plans, the first registry request after the server starts builds a bounded in-memory index from \code{model_record.json}, \code{run_nodes.jsonl}, and \code{run_edges.jsonl}.
Subsequent page, family, and deep-link requests reuse that index instead of reading and parsing the complete files again.
Before reuse, the server compares each source file's absolute path, size, and modification time; a changed source is reparsed automatically, while an unchanged source continues to use the cached index.
Only the active runtime map and resolved plan are retained, so selecting another simulator or policy replaces the previous cache entry.
The first request after a server restart or source-file change is therefore a cold load and can take longer than later requests.

\section*{Paper Evaluation View}

The dedicated \code{/evaluation} route is separate from the generation-pool view at \code{/}.
It visualizes only the exact released samples used for paper evaluation and shows the ground-truth label by default.
The web request path does not run simulations or agents.
Navigation links in both views switch between the generation pool and paper evaluation without changing either data source.

The paper-evaluation controls select:

\begin{itemize}
\item environment: \code{BallDrop}, \code{BounceBall}, or \code{MassSlide};
\item collection: the released task-test samples or the balanced 13-per-class batch;
\item seed 0--4 or all seeds;
\item ground-truth class;
\item clean, low-noise, or high-noise trajectories and a deterministic noise seed;
\item independent clean-baseline, time-zero-baseline, and
alternative-root-cause overlays.
\end{itemize}

The task-test view contains ten released test samples per seed and exactly 50 unique samples per environment across seeds 0--4.
The balanced view contains 13 samples per class per seed: 65 for BallDrop, 78 for BounceBall, and 65 for MassSlide.
The application verifies that the released 0-shot and 3-shot IDs are identical and presents them as one shared collection.
The 1-shot \code{fern_01234} rows are never loaded into this view.

Each sample row shows its ID, ground-truth class, and seed membership.
The detail area shows the released trajectory, intervention metadata, class counts, previous/next navigation, and a shareable URL.
When the paired baseline Parquet is present in the released bundle, the user can select \code{Overlay clean baseline}, and noise diagnostics can report global SNR.
If that Parquet was not released, only the baseline overlay and baseline-dependent SNR are disabled; clean and noisy sample visualization remain available.
The separate \code{Overlay time-0 baseline} checkbox plots the same selected
intervention magnitude applied cleanly from simulation time zero.
The two checkboxes are independent, so the selected released episode, clean
parent baseline, and time-zero counterpart can appear together.
The \code{Overlay alternative-root-cause probe} checkbox adds the best-fitting
trajectory obtained by excluding the ground-truth changed parameter,
optimizing every other exposed parameter over its configured legal interval,
and applying the winning change at the released episode's original
intervention time.
The viewer displays the winning alternative parameter, value, and normalized
optimization loss.
No-intervention samples do not have this probe because there is no
ground-truth changed parameter to exclude.
All three checkboxes are independent.
Observation noise applies only to the selected released episode; every probe
overlay remains clean.

The time-zero trajectories are auxiliary viewer artifacts, not members of the
released evaluation set.
They are deterministically regenerated from the validated
\code{model_run_specs.json} mappings with MATLAB R2025a Update 1, without rerunning any
benchmark agent.
The materializer validates the stored parent, child, and time-zero hashes,
checks every generated Parquet against the simulator signal schema and time
grid, and records source hashes, recipes, MATLAB versions, file hashes, and row
counts in \code{time0_auxiliary_manifest.json}.
Across the exact task-test and balanced released selections, there are 89
BallDrop, 106 BounceBall, and 95 MassSlide time-zero trajectories (290 total).
They can be reproduced with:

\begin{DocCode}
SIMSCAPE_CUSTOM_PATH=/path/to/tsENV/simscape_custom \
  python workflows/simulate/materialize_evaluation_time0_baselines.py \
  --balanced-manifest /path/to/balanced_heldout_evaluation.json \
  --output-root /path/to/evaluation_time0_aux
\end{DocCode}

The optimized alternative-root-cause trajectories use the same 290 released
intervention episodes.
They are stored in another isolated auxiliary artifact and do not change
evaluation membership or agent outputs.
The materializer uses the exact released clean episode as its target, searches
all non-ground-truth parameters with a nine-point coarse grid followed by
\code{fminbnd} refinement in an explicitly verified eight-worker MATLAB
\code{Processes} pool, verifies every optimized value against its legal
interval, and validates the resulting Parquet against the simulator's signal
schema and time grid.
Its manifest records every candidate report, winning parameter and value,
loss, source hashes, recipe, MATLAB version, and Parquet checksum.
It can be reproduced with:

\begin{DocCode}
SIMSCAPE_CUSTOM_PATH=/path/to/tsENV/simscape_custom \
  python workflows/simulate/materialize_evaluation_alternative_probes.py \
  --balanced-manifest /path/to/balanced_heldout_evaluation.json \
  --time0-root /path/to/evaluation_time0_aux \
  --matlab-workers 8 \
  --output-root /path/to/evaluation_alternative_probe_aux
\end{DocCode}

Examples:

\begin{DocCode}
http://localhost:3001/evaluation?environment=BallDrop&collection=task_test&seed=0
http://localhost:3001/evaluation?environment=BallDrop&collection=balanced&seed=0
http://localhost:3001/evaluation?environment=BallDrop&collection=task_test&seed=0&sample=<RUN_ID>
\end{DocCode}

Before returning any samples, the evaluation registry verifies the source hashes recorded by the balanced manifest, the released labels and parameter hashes, per-seed sizes, 13-per-class counts, sample-to-baseline mappings, released Parquet presence, held-out membership, and 0-shot/3-shot identity.
A mismatch produces an explicit error instead of reconstructing or guessing a selection.
Trajectory requests from this page use \code{source=evaluation}, which prevents a same-named generation artifact from replacing a released paper input.
Time-zero overlay requests use \code{source=evaluation_time0}; this source is
strictly isolated from both released and generation artifacts and accepts only
the clean noise profile.
Alternative-root-cause requests use
\code{source=evaluation_alternative}; this source is likewise isolated and
clean-only.
The Python noise loader accepts both the current \code{baseline=} interface and the legacy \code{ref=} interface.



\section*{What It Shows}

The sidebar lists allowed \name{} \termSimulators that exist under \code{models/simulink/}.
\termSimulators with missing required files remain listed but show a warning badge.
The header shows the selected \termSimulator, selected \code{policy_id}, active runs folder name, and a reload button. 
The reload button reloads the visible page of families and any expanded rows.
This revalidates the source-file metadata: unchanged files reuse the in-memory index, while changed files trigger a cold rebuild.

The main table is organized by \code{family_id} from the resolved run graph.
Each family can be expanded to show its baseline run, intervention runs, and matched time-zero/static-change probe runs.

Table columns include:

\begin{itemize}
\item status from \code{model_record.json};
\item cheap-filter status from \code{cheap_filter_metrics.jsonl};
\item \code{run_id};
\item \code{family_id};
\item \code{source};
\item \code{validation_profile};
\item \code{recipe_hash};
\item exposed parameter and initial-state values in experiment-config order;
\item actions for plotting and expansion.
\end{itemize}

Expanded intervention rows show changed parameter, direction, new value, and intervention time from \code{run_nodes.jsonl} and \code{run_edges.jsonl}.
For production policies with surrogate scoring, the table also shows \code{surrogate_id}, \code{noise_level}, \code{predicted_label}, \code{true_label_confidence}, \code{confidence_margin}, and \code{surrogate_filter_pass}.

\section*{Plot Area}

The plot area renders selected runs.
Plot data can be shown with no noise or with a \termSimulator-provided low or high noise profile when \code{noise_adder.py} exists.

Controls include:

\begin{itemize}
\item run-selection eye buttons in the table;
\item noise profile: \code{none}, \code{low}, or \code{high};
\item deterministic noise seed;
\item observed signal selection.
\end{itemize}

If two runs are selected and noise is enabled, the UI may show per-signal SNR values using the \termSimulator noise-adder utilities.

\section*{Deep Links}
Deep links can preselect a \termSimulator, policy, run, and optional comparator.
When a deep link targets a run, the explorer resolves the run and fetches only the containing family needed to display it, rather than loading the full policy registry.
Supported query parameters include:

\begin{DocCode}
model, policy, run, run2, compare_run, compare, noise_profile, noise_seed
\end{DocCode}

The \code{compare} parameter may be \code{baseline}, \code{time0}, \code{sibling}, or \code{none}.
A plain \code{?run=<RUN_ID>} link opens a single-run view unless an explicit comparator is provided.

Shortcut:

\begin{DocCode}
http://localhost:3001/<run_id>
\end{DocCode}

opens the matching run directly if it can be located.

\section*{Endpoints Used}

\begin{itemize}
\item \code{GET /api/models}: read allowed \termSimulators and validation warnings.
\item \code{GET /api/policies?model=<simulator_id>}: list available resolved plans under \code{plans/}.
\item \code{GET /api/registry?model=<simulator_id>&policy=<POLICY_ID>}: load the first page of family summaries for the selected policy.
\item \code{GET /api/registry?model=<simulator_id>&policy=<POLICY_ID>&page=<N>&page_size=<N>}: load a specific page of family summaries.
\item \code{GET /api/registry?model=<simulator_id>&policy=<POLICY_ID>&eligible_only=true}: load only families with passing cheap-filter results and, when family details are requested, include only their passing intervention runs.
\item \code{GET /api/registry?model=<simulator_id>&policy=<POLICY_ID>&family_id=<FAMILY_ID>}: load full details for one family, including baseline, intervention, and probe runs.
\item \code{GET /api/registry?model=<simulator_id>&policy=<POLICY_ID>&run=<RUN_ID>}: resolve a run within a policy and load the containing family for deep-link display.
\item \code{GET /api/registry?model=<simulator_id>&policy=<POLICY_ID>&mode=full}: explicitly load the full registry for debugging or compatibility.
\item \code{GET /api/distribution?model=<simulator_id>}: load \code{experiment_config.json} and generated metadata used for legal domains and sampling details.
\item \code{GET /api/locate-run?run=<RUN_ID>}: resolve a deep-linked run identifier to \termSimulator and policy.
\item \code{GET /api/data?model=<simulator_id>&run=<RUN_ID>&max_rows=<N>&noise_profile=<PROFILE>&noise_seed=<SEED>}: load a run artifact for plotting.
\item \code{GET /api/evaluation-registry?environment=<ENVIRONMENT>&collection=<task_test|balanced>&seed=<0--4|all>}: validate and load the exact released paper-evaluation sample registry; an optional \code{class=<LABEL>} filters the returned rows.
\item \code{GET /api/data?model=<ENVIRONMENT>&run=<RUN_ID>&source=evaluation}: load only the released paper-evaluation Parquet for plotting.
\item \code{GET /api/data?model=<ENVIRONMENT>&run=<RUN_ID>&source=evaluation_time0}: load only a clean, regenerated time-zero auxiliary Parquet for plotting.
\item \code{GET /api/data?model=<ENVIRONMENT>&run=<RUN_ID>&source=evaluation_alternative}: load only a clean, optimized alternative-root-cause auxiliary Parquet for plotting.
\item \code{GET /api/simulate?model=<simulator_id>&policy=<POLICY_ID>}: read saved simulation status information when present.
\end{itemize}

\section*{Files Read}

\begin{itemize}
\item \code{models/simulink/<simulator_id>/plans/<policy_id>/run_nodes.jsonl}
\item \code{models/simulink/<simulator_id>/plans/<policy_id>/run_edges.jsonl}
\item \code{models/simulink/<simulator_id>/plans/<policy_id>/generation_report.json}
\item \code{models/simulink/<simulator_id>/plans/<policy_id>/validation_report.json}
\item \code{models/simulink/<simulator_id>/runs/model_record.json}
\item \code{models/simulink/<simulator_id>/metrics/<policy_id>/cheap_filter_metrics.jsonl}
\item \code{models/simulink/<simulator_id>/metrics/<policy_id>/surrogate_scores/<surrogate_id>.jsonl}
\item \code{models/simulink/<simulator_id>/surrogates/<surrogate_id>/thresholds.json}
\item \code{models/simulink/<simulator_id>/generated/metadata.json}
\item \code{models/simulink/<simulator_id>/description_levels.json}
\item \code{models/simulink/<simulator_id>/experiment_config.json}
\item \code{models/simulink/<simulator_id>/runs/<run_id>/data.parquet}
\item \code{TraceBench_questions/<simulator_id>/dataframes/<run_id>.parquet}, as fallback for exported bundles
\item \code{models/simulink/<simulator_id>/noise_adder.py}
\item \code{<evaluation_root>/<environment>/sample_manifest.json}
\item \code{<evaluation_root>/<environment>/model_record.json}
\item \code{<evaluation_root>/<environment>/questions.json}
\item \code{<evaluation_root>/<environment>/noise_adder.py}
\item \code{<evaluation_root>/<environment>/dataframes/<run_id>.parquet}
\item the \code{balanced_heldout_evaluation.json} file named by \code{WEB_MODEL_EXPLORER_PAPER_EVALUATION_MANIFEST}
\item \code{<evaluation_time0_root>/<environment>/time0_baselines/<time0_run_id>.parquet}
\item \code{<evaluation_time0_root>/time0_auxiliary_manifest.json}
\item \code{<evaluation_alternative_probe_root>/<environment>/alternative_root_cause/<probe_run_id>.parquet}
\item \code{<evaluation_alternative_probe_root>/<environment>/optimisation_reports/<sample_run_id>.json}
\item \code{<evaluation_alternative_probe_root>/alternative_probe_auxiliary_manifest.json}
\item \code{models/simulink/<simulator_id>/model_run_specs.json}
\end{itemize}


\section*{Human Study}
The Human Study app checks whether exported questions are feasible for humans.
It reads from \code{TraceBench_questions/} and can be started with:

\begin{DocCode}
bash web_human_study/start.sh [--port PORT]
\end{DocCode}

The app visualizes time series with Plotly and shows the multiple-choice labels from \code{questions.json}.
When a user selects an option, the correct label is shown, and the result is saved under \code{human_feedback/<name>.json}.
There is also a button to mark a question as problematic.

Shortcut forms include:

\begin{itemize}
\item \code{http://localhost:5173/<run_id>}: open a sample by \termRunID.
\item \code{http://localhost:5173/<question_id>}: open a question.
\item \code{http://localhost:5173/<question_id>-<run_id>}: open a specific run within a question.
\end{itemize}

\end{document}
